\documentclass[10pt,a4paper]{amsart}
\usepackage[margin=19mm]{geometry}
\usepackage{amsmath,amssymb,mathtools}
\usepackage[scr=rsfs,cal=euler]{mathalfa}
\usepackage{xfrac}
\usepackage[unicode,hidelinks,bookmarks=false]{hyperref}
\newcommand{\ie}{i.e.\ }
\newcommand{\cf}{cf.\ }
\newcommand{\resp}{respectively\ }
\newcommand{\Subsection}{\subsection}
\providecommand{\nobreakdash}{\nobreak\hbox{-}}
\setlength{\emergencystretch}{2em}
\multlinegap0pt
\usepackage{tikz}
\usepackage{amssymb}
\usepackage[shortlabels]{enumitem}
\usepackage{enumerate}
\usepackage[normalem]{ulem}
\usepackage{algorithm}\usepackage{algpseudocode}
\usepackage{float}
\newtheorem{corollary}{Corollary}
\newtheorem{theorem}{Theorem}
\newcommand{\es}{\alpha_{\rm es}}
\title{On externally supported independence number of graphs}
\author{Dragana Bo\v zovi\'{c}\textsuperscript{a,b} \and Iztok Peterin \textsuperscript{a.c} \and Adriana Roux \textsuperscript{d} \and Aleksandra Tepeh\textsuperscript{a,b,c}}
\date{Source: arXiv:2606.22972v1 (CC BY 4.0)}
\begin{document}
\maketitle

\noindent\emph{Source and licence.} On externally supported independence number of graphs. Version arXiv:2606.22972v1 (CC BY 4.0), distributed by arXiv under the Creative Commons Attribution 4.0 International licence, http://creativecommons.org/licenses/by/4.0/, as recorded in the arXiv licence metadata for this exact version and shown on the abstract page https://arxiv.org/abs/2606.22972v1. Full licence notice: \texttt{licenses/arxiv-2606.22972-CC-BY-4.0.txt}.\par\smallskip

\noindent\textit{Editorial scope.} Counted unit: the article's theorem tree-lower (source0.tex lines 784--787). The article's complete original proof of it is delivered with it (source0.tex lines 789--833, one \texttt{proof} environment of 45 lines). No mathematics is written, repaired or re-derived: the statement and the proof are the article's own lines, in the article's own notation, and no retained line is substituted or re-flowed.  Retained uncounted as the source-local context the proof needs: lines 432--436.  Nothing was omitted from any retained span, so the package carries no omission record.  No retained line contains a citation, so the wrapper carries no bibliography and no bibliography entry.  No retained line contains a figure environment, an image-inclusion command or drawing code, so no image is delivered and the package declares no asset dependency.  Editorial formatting only, in addition: an AMS \texttt{amsart} preamble that restates the article's own declarations and macros used by the retained lines, namely the environments \texttt{corollary}, \texttt{theorem} on separate counters, together with the standard packages those lines need.  The retained environments print in this document as Corollary 1, Theorem 1 in order of appearance, whereas the article prints the same items with the numbering of its own full document.  The complete native source of the article as distributed in the arXiv e-print for this exact version is bundled unchanged as \texttt{sources/203396/source0.tex}.
\begin{corollary}
\label{notESI1}
If $B$ is an ESI-set of a graph $G$, then $B$ contains neither a universal vertex, nor a
support vertex, nor a twin, nor a support of a twin.
\end{corollary}

\begin{theorem}
\label{tree-lower}
If $T$ is a tree, then $\es(T)\geq \ell(T)-s(T)$, where equality holds if and only if $V(T)=S(T)\cup L(T)$ and $s(v) \leq \ell(v)-1$ for every $v\in S(T)$.
\end{theorem}

\begin{proof}
Let $T$ be a tree and let $B$ be its canonical ESI-set.  So,  $|B|=\ell(T)-s(T)$ and the lower bound $\es(T)\geq |B|=\ell(T)-s(T)$ follows. 
We next characterize when the equality holds.

We first prove that if $V(T)-(S(T)\cup L(T))\neq \emptyset$, then $\es(T)>\ell(T)-s(T)$.

Let $v\in V(T)-(S(T)\cup L(T))$ and let $z$ be a leaf at minimum distance from $v$ with $v=v_0,v_1,\dots,v_k=z$ the unique $v$--$z$ path. Then $v_{k-1}\in S(T)$. Since $v$ is not a support vertex, we have $k\ge 2$, and hence $u=v_{k-2}$ is a non-leaf vertex with $\ell(u)=0$ and $w=v_{k-1}\in S(u)$. 
%(If $k=2$, then $u=v$.)
In a canonical ESI-set $B$, the chosen leaf $t_w$ satisfies $t_w\notin B$. Define $B^\star=B\cup\{t_w\}$, which is clearly an independent set. Then we have $|B^\star|=|B|+1$. Let $D^\star=V(T)-N[B^\star]$. Since $u$ is not adjacent to any leaf, it is not adjacent to any vertex of $B^\star$, and thus $u\in D^\star$.
We now verify the ES-condition for $B^\star$. Every vertex of $N(B)$ keeps an external neighbor in $D^\star$ as in the canonical construction. Moreover, the only vertex that may be newly added to $N(B^\star)$ is $w$, 
%(which happens in the case when $l(w)=1$),
which has a neighbor in $D^\star$, namely $u$. Hence $B^\star$ is an ESI-set and therefore $\es(T)\ge |B^\star|=|B|+1>|B|=\ell(T)-s(T)$.

Now we prove that the existence of a support vertex $v$ with $s(v)>\ell(v)-1$ forces the strict inequality $\es(T)>\ell(T)-s(T)$. Let $v\in S(T)$ with $s(v)>\ell(v)-1$ and perform a switch at $v$ on the canonical pair $(B,D)$ to obtain $(B_v,D_v)$.
Since $B\cap L(v)=L(v)-\{t_v\}$, we remove $\ell(v)-1$ leaves from $B$ and add $s(v)$
leaves (the vertices $t_w$ for every $w\in S(v)$) to obtain $B_v$. Therefore
$|B_v|=|B|-(\ell(v)-1)+s(v) > |B|$, since $s(v)>\ell(v)-1$, 
and hence $\es(T)\ge |B_v|>|B|=\ell(T)-s(T)$. 

Conversely, assume that $V(T)=S(T)\cup L(T)$ and $s(v)\le \ell(v)-1$ holds for every $v\in S(T)$.
Suppose for a contradiction that $\es(T)>\ell(T)-s(T)$, and let $B$ be an $\es(T)$-set with
$D=V(T)-N[B]$. Among all $\es(T)$-sets choose $B$ so that $|S(T)\cap D|$ is minimum.

If $S(T)\cap D\neq\emptyset$, pick $x\in S(T)\cap D$ and perform a reverse switch at $x$,
obtaining an ESI-set $B^{x}$. By construction,
$|B^{x}|=|B|-s(x)+(\ell(x)-1)\ \ge\ |B|$. Since $|B|=\es(T)$, we have $|B^{x}|=|B|$, while the reverse
switch removes $x$ from $D\cap S(T)$, contradicting the choice of $B$. Hence
$S(T)\cap D=\emptyset$.

Since $V(T)=S(T)\cup L(T)$ and $S(T)\cap D=\emptyset$, it follows that
$D\subseteq L(T)$, i.e.~every vertex of $D$ is a leaf.

We may further assume that $B\subseteq L(T)$. Indeed, if $x\in B$ were a non-leaf vertex, then $x\in S(T)$,
contradicting Corollary~\ref{notESI1}, which states that an ESI-set contains no support vertex.

Now fix a support vertex $w\in S(T)$. If $B$ contains all leaves of $L(w)$, then $w\in N(B)$,
and the ES-condition requires $w$ to have a neighbor in $D$, a contradiction.

Therefore $|B\cap L(w)|\le \ell(w)-1$ for every $w\in S(T)$, and hence
$$|B|
=\sum_{w\in S(T)}|B\cap L(w)|
\le \sum_{w\in S(T)}(\ell(w)-1)
=\ell(T)-s(T),$$
contradicting $|B|=\es(T)>\ell(T)-s(T)$. 
\end{proof}

\end{document}
